% HPO-Adapt tool paper (~4 pages + references).
% Compile: pdflatex paper && bibtex paper && pdflatex paper && pdflatex paper
%
% NOTE ON NUMBERS: every numeric claim in the Evaluation section is taken
% verbatim from experiments/results/full/figs/*.csv (see paper/FIGURES.md).
% Regenerate with: python experiments/run_all.py --tag full && python
% experiments/make_figures.py --tag full

\documentclass[10pt]{article}
\usepackage[margin=0.85in]{geometry}
\usepackage[T1]{fontenc}
\usepackage[utf8]{inputenc}
\usepackage{microtype}
\usepackage{booktabs}
\usepackage{amsmath}
\usepackage{amssymb}
\usepackage{graphicx}
\usepackage{xcolor}
\usepackage[hidelinks]{hyperref}
\usepackage{caption}
\captionsetup{font=small,labelfont=bf}
\linespread{0.97}

\newcommand{\toolname}{HPO-Adapt}
\newcommand{\todofig}[1]{\fbox{\parbox[c][0.55cm][c]{0.94\linewidth}{\small\textbf{[FIGURE:]} #1}}}

\title{\toolname{}: A Diagnostic-Driven Toolkit for Adaptive\\ Hyperparameter-Optimization Setup Control}
\author{York Wang\thanks{Course project; the toolkit, experiments and this paper were produced as
part of a graduate AutoML course project. Code: \url{https://github.com/yorkwang/hpo-adapt}
(replace with the final repository URL).}}
\date{}

\begin{document}
\maketitle

\begin{abstract}
The performance of hyperparameter optimization (HPO) is largely determined
before the first trial runs: the initial \emph{setup}---search space,
optimizer, budget---fixes what can be found at all. Analysis frameworks
such as CAVE and DeepCAVE expose what went wrong \emph{afterwards} but
leave the repair to the user, while LLM-agent approaches tend to replace
classical HPO rather than repair its setup. We present \toolname{}, a
Python toolkit that closes this loop automatically: a deterministic
\emph{Diagnostic Engine} computes a structured state from the trial
trajectory, a tabular Q-learning \emph{Controller} selects one of seven
high-level interventions, a \emph{Revised Setup Agent} executes it as a
validated new setup, and an \emph{Agent Memory} retains experiences across
episodes. Every agent has a deterministic offline backend, so the full
loop runs without any API key and is exactly reproducible. On two
benchmark tasks whose initial setup is misspecified, guided adaptation
reaches substantially better optima than equally-budgeted static HPO
(1.000 vs.\ 0.755; 0.498 vs.\ 0.896, approaching the global minimum),
while on a real RandomForest task it matches strong static baselines
within noise; the controller keeps improving across episodes, and the
observe--decide layer costs 1.4\% of wall-clock on real training
workloads.
\end{abstract}

\section{Introduction}
\label{sec:intro}

Hyperparameter optimization automates the search for good model
configurations, but it silently inherits every flaw of its initial setup:
a search space that is too narrow truncates the optimum, an unsuited
optimizer or budget wastes the search \cite{feurer2019hyperparameter}.
Practitioners usually discover such flaws only after the budget is gone,
by inspecting the HPO trajectory with tools such as CAVE
\cite{eggensperger2019cave} or DeepCAVE \cite{segel2025deepcave}. These
frameworks make HPO transparent, but \emph{repairing} the setup remains a
manual, experience-driven step. A second line of work assigns the setup
to LLM agents \cite{liu2025agenthpo,trirat2025automlagent}; recent
evidence indicates that LLMs do not outperform classical HPO algorithms
as searchers \cite{ferreira2026llms}, and such agents remain bound by the
bounds and interfaces they are given \cite{tornede2024automl}. Neither
line answers the question that matters during a running search:
\emph{given the trajectory so far, which concrete change to the current
setup is most promising right now?}

\toolname{} answers it with a closed loop and a strict separation of
concerns. \textbf{Observe}: a deterministic \emph{Diagnostic Engine}
reduces the cumulative trial trajectory to a structured
\texttt{DiagnosticState} (importance, boundary concentration, good
regions, diversity, coverage, budget), with explicit
\texttt{VALID}/\texttt{INSUFFICIENT\_DATA} statuses instead of fabricated
numbers. \textbf{Decide}: a tabular Q-learning \emph{Controller}
\cite{sutton2018reinforcement} maps that state to one of seven high-level
interventions. \textbf{Execute}: a \emph{Revised Setup Agent} turns the
intervention into a new setup that must pass a tool validator before any
HPO run; an \emph{Agent Memory} (SQLite) stores completed adaptation
experiences. Both agents also ship deterministic heuristic backends, so
the system runs offline, is seeded end-to-end, and every number in this
paper regenerates with one command. Contributions: \textbf{(i)}~a modular
open-source toolkit that turns CAVE-style trajectory diagnosis into
\emph{acted-on} setup adaptation rather than after-the-fact plots;
\textbf{(ii)}~a design that keeps metric computation free of LLMs and the
agent role strictly at validated setup construction; \textbf{(iii)}~126
unit tests and one-command reproducible experiments; and \textbf{(iv)}~an
empirical study showing that, on misspecified setups, guided adaptation
clearly beats equally-budgeted static HPO while unguided intervention does
not, that the controller improves across episodes, and that the
observe--decide layer costs 1.4\% of wall-clock on real training
workloads.

\section{Related Work}
\label{sec:related}

\paragraph{HPO and its setup problem.}
Classical HPO samples configurations from a fixed search space
\cite{feurer2019hyperparameter}; Auto-WEKA already framed the joint choice
of model and hyperparameters as a search problem
\cite{thornton2013autoweka}. Practical HPO relies on model-based tools
such as SMAC3 \cite{lindauer2019smac3} and Optuna
\cite{akiba2019optuna}; \toolname{} builds on Optuna and targets the
situation these tools leave to the user: the setup was chosen once and
may be wrong.

\paragraph{Trajectory analysis.}
CAVE \cite{eggensperger2019cave} and DeepCAVE \cite{segel2025deepcave}
assess finished HPO runs (importance, configuration footprints): they
inform humans but do not feed back into the running search. \toolname{}
reuses their core indicators---fANOVA-style importance and
local-perturbation importance---but recomputes them \emph{during} the
search, on the cumulative trajectory, as one structured state that a
controller can act on.

\paragraph{LLM agents for AutoML.}
AgentHPO \cite{liu2025agenthpo} and AutoML-Agent \cite{trirat2025automlagent}
delegate setup and even search to LLM agents; such agents remain bound by
externally imposed bounds and interfaces \cite{tornede2024automl}, and
controlled studies find classical HPO still stronger as the searcher
\cite{ferreira2026llms}. \toolname{} keeps classical HPO as the searcher
and assigns agents only the setup, with every output schema-validated---a
design that also makes the system runnable fully offline through
deterministic heuristic backends.

\section{The \toolname{} Toolkit}
\label{sec:system}

\subsection{Closed loop and data objects}
\label{sec:loop}

\todofig{Fig.~1: closed loop
Task $\to$ Initial Setup $\to$ $S_0$ $\to$ HPO (Optuna) $\to$ TrialLogs
$\to$ Diagnostic Engine $\to$ RL Controller $\to$ Revised Setup $\to$
$S_{t+1}$, with the Tool Validator gate and SQLite memory. Full spec:
paper/FIGURES.md Fig.~1.}

The system is organized around six pydantic schemas: \texttt{TaskSpec}
(objective, constraints, and the \emph{allowed} hyperparameter metadata),
\texttt{Setup} (search space, optimizer, budget, evaluation),
\texttt{TrialLog}, \texttt{DiagnosticState}, \texttt{RLAction}
($A_t=(\text{ActionType},\text{Target},\text{Parameters})$), and
\texttt{Experience}. Every setup must pass \texttt{validate\_setup} before
the HPO runner accepts it: the validator rejects invented hyperparameters,
out-of-domain ranges, type or direction mismatches, unknown optimizers,
and budget violations.

\subsection{Diagnostic Engine (Observe)}
\label{sec:diagnostics}

The engine is pure, deterministic Python; no LLM sees raw logs. It
analyzes the \emph{cumulative} trajectory of the current episode and
returns a state with three sections. \emph{Optimization status}: best
objective $B_t=\max_{i\le t}\tilde y_i$, recent improvement over an
analysis window $w$, $RI_t=(B_t-B_{t-w})/(|B_{t-w}|+\epsilon)$,
stagnation $\mathbb{1}(RI_t<\epsilon_{\mathrm{stag}})$, and marginal gain
per trial/second, computed on internally maximized values
$\tilde y_i=y_i$ (maximize) or $-y_i$ (minimize). \emph{Search-space
status}: per hyperparameter, (i)~global importance from an fANOVA-style
variance decomposition over a random-forest surrogate plus local
importance from leave-one-out prediction perturbation at the best
configurations \cite{hutter2014efficient}; (ii)~boundary concentration of
all trials and of the top-$q\%$ trials (per upper/lower edge)---the key
truncation signal; (iii)~the \emph{good region} (mean and $[10\%,90\%]$
quantile interval, width ratio vs.\ the current range; category
frequencies for categoricals); (iv)~configuration diversity (mixed-type
pairwise distance of the recent window vs.\ the full trajectory) and
(v)~log-space-aware per-dimension coverage.
\emph{Budget status}: consumed/remaining share of the global budget.
Every metric carries a status: below the minimum trial counts (5/20/30
for basic/stagnation/importance metrics) the engine reports
\texttt{INSUFFICIENT\_DATA} instead of a value, and hyperparameter-targeted
actions stay unavailable until importance becomes \texttt{VALID}---the
system never fabricates evidence.

\subsection{RL Controller (Decide)}
\label{sec:rl}

The global state discretizes stagnation, recent improvement, diversity,
coverage, budget remaining, and marginal gain into low/medium/high; up to
two candidate hyperparameters (ranked by
$\alpha\,\mathrm{Imp}_j+\beta\,BC^{\mathrm{good}}_j+
\gamma\,(1-\mathrm{WidthRatio}_j)$, $\alpha=.5,\beta=.3,\gamma=.2$) add
per-candidate context. Actions are the seven interventions in
Table~\ref{tab:actions}. Selection is $\epsilon$-greedy with decay; ties
among equally-good actions are broken by a rule engine that encodes seven
diagnostic patterns (e.g., \emph{stagnation + high good-boundary
concentration + high importance} $\Rightarrow$ \textsc{Expand}; \emph{low
importance for multiple rounds} $\Rightarrow$ \textsc{Freeze}), so the
rules act as an interpretable prior while Q-values refine it from
experience. \textsc{Stop} is evidence-gated: selectable only under observed
stagnation with at least half the budget consumed, preventing unguided
premature termination. The reward is
$r=\mathrm{clip}(\Delta-\lambda C-r_{\mathrm{invalid}},\pm r_{\max})$ with
$\Delta=(B_{t+1}-B_t)/(|B_t|+\epsilon)$ the relative improvement caused by
the intervention, $C$ the budget share it consumed, and a fixed penalty
for structurally invalid actions (which never consume budget and are
retried within the round).

\begin{center}
\small
\captionof{table}{The seven setup interventions ($A_t=$ action, target, parameters).}
\label{tab:actions}
\begin{tabular}{@{}ll@{}}
\toprule
Action (target) & Semantics \\
\midrule
\textsc{ExpandSearchSpace} (numeric hp) & widen range toward the boundary indicated by good trials \\
\textsc{ShiftSearchRegion} (numeric hp) & translate range toward the observed good region \\
\textsc{ShrinkSearchSpace} (numeric hp) & focus range on the good region (with margin) \\
\textsc{FreezeHyperparameter} (any hp) & fix a persistently low-importance parameter \\
\textsc{IncreaseExploration} (optimizer) & raise exploration / switch optimizer \\
\textsc{IncreaseBudget} (budget) & enlarge the next phase's trial budget (capped) \\
\textsc{Stop} (---) & terminate the episode (evidence-gated) \\
\bottomrule
\end{tabular}
\end{center}

\subsection{Setup Agents (Execute) and Memory}
\label{sec:agents}

The \emph{Initial Setup Agent} constructs $S_0$ from task metadata,
constraints, and memory context; the \emph{Revised Setup Agent} applies the
selected intervention \emph{conservatively}: only the targeted component
changes, expansions are clipped to the task's metadata domain, shifts
preserve width, shrinks never fall below a minimum relative width, and
freezes prefer the task default. Both agents have a deterministic
heuristic backend (default; used for all experiments) and an LLM backend
(any OpenAI-compatible endpoint) with identical interfaces; every output
passes the same validator, so nothing violating the schema or task
constraints is ever executed. Phases are \emph{warm-started}: each new
Optuna study is seeded with all past completed trials compatible with the
current search space, so a revised setup inherits every evaluation instead
of restarting the optimizer's random startup; only newly evaluated trials
consume budget, keeping static and adaptive runs comparable.

The \emph{Agent Memory} is a SQLite store with two tables:
\texttt{experiences} (one row per completed adaptation step: task
signature, setup before/after, diagnostic state, diagnosis, action,
outcome, reward) and \texttt{semantic\_rules}, which aggregates outcomes
into rule-level success rates. By default only successful interventions and
instructive failures are persisted. Memory answers ``what happened in
similar situations before?'' while the Q-table answers ``what pays off in
this state?''---complementary roles.

\subsection{Using the toolkit}

\toolname{} installs with \texttt{pip install -e .} and offers a CLI
(\texttt{hpo-adapt --task synthetic\_mlp --max-rounds 4 --seed 3}), a
Python API (\texttt{AdaptationPipeline(task, PipelineConfig()).run()}), and
plug-in points for custom objectives: register the objective function,
describe the task's hyperparameter metadata in a \texttt{TaskSpec}, and run
the pipeline. The shipped tasks cover a synthetic MLP-style benchmark
whose optimum lies outside a default-derived learning-rate range, the
Branin function (minimize), and a real RandomForest/3-fold-CV task on
sklearn digits.

\section{Evaluation}
\label{sec:eval}

\subsection{Setup}

We compare, under an \emph{identical} total budget of 120 evaluations per
run and 10 seeds each: static baselines that spend the whole budget in one
phase (\texttt{static-random}, \texttt{static-tpe}) and four adaptive
variants (phases of 40 trials, at most three phases):
\texttt{adapt-random-action} (unguided control), \texttt{adapt-rule-only}
(always the rule engine's recommendation), \texttt{adapt-rl-default}
(spec-default $\epsilon$-greedy $1.0\!\to\!0.05$), and
\texttt{adapt-rl-guided} ($\epsilon_0=0.2$). A second experiment shares
Q-table and memory across 20 consecutive episodes per (task, seed). All
runs use the deterministic heuristic agent backends; every number below
regenerates with \texttt{run\_all.py} + \texttt{make\_figures.py}.

\subsection{Main results}

\begin{center}
\small
\captionof{table}{Final best objective (mean$\pm$std, 10 seeds, raw space;
Branin minimized; best per task in \textbf{bold}).}
\label{tab:main}
\begin{tabular}{@{}lccc@{}}
\toprule
 & synthetic\_mlp ($\uparrow$) & Branin ($\downarrow$) & rf\_digits ($\uparrow$) \\
\midrule
static-random   & 0.657$\pm$0.039 & 1.759$\pm$0.426 & 0.9766$\pm$0.0008 \\
static-tpe      & 0.755$\pm$0.003 & 0.896$\pm$0.010 & \textbf{0.9772$\pm$0.0004} \\
adapt-random-action & 0.873$\pm$0.134 & 0.776$\pm$0.258 & 0.9767$\pm$0.0008 \\
adapt-rule-only & \textbf{1.000$\pm$0.000} & \textbf{0.498$\pm$0.198} & 0.9771$\pm$0.0005 \\
adapt-rl-default & 0.873$\pm$0.134 & 0.776$\pm$0.258 & 0.9767$\pm$0.0008 \\
adapt-rl-guided & \textbf{1.000$\pm$0.000} & 0.596$\pm$0.251 & 0.9771$\pm$0.0005 \\
\bottomrule
\end{tabular}
\end{center}

\todofig{Fig.~2: mean best-so-far vs.\ trials, 6 methods $\times$ 3 tasks;
dashed = static. Data: \texttt{figs/anytime.csv} (FIGURES.md Fig.~2).}

On the synthetic task the picture is unambiguous: guided adaptation
reaches the achievable optimum (1.000$\pm$0.000, i.e., every seed recovers
learning rate $\approx 0.42$--$0.45$), while static TPE tops out at
0.755$\pm$0.003---a 32\% relative gap. On Branin, \texttt{adapt-rule-only}
reaches 0.498$\pm$0.198 against static TPE's 0.896$\pm$0.010, approaching
the global minimum 0.398. On the real task all methods land within 0.0006
of each other (0.9766--0.9772): the metadata-driven initial setup is
already near the CV-accuracy plateau, so adaptation neither helps nor
hurts---the toolkit's value there is observability.

\subsection{Ablations: guidance matters}

Comparing the three adaptation policies isolates the contribution of
guidance. On synthetic\_mlp, \texttt{adapt-rule-only} and
\texttt{adapt-rl-guided} achieve 1.000$\pm$0.000 (worst seed = best seed),
whereas \texttt{adapt-random-action} manages only 0.873$\pm$0.134---a
six-fold larger spread. On Branin the ordering is the same:
0.498$\pm$0.198 and 0.596$\pm$0.251 for the guided variants against
0.776$\pm$0.258 for the random control. Intervening \emph{without}
evidence buys little and costs variance. Under spec-default exploration,
\texttt{adapt-rl-default} behaves exactly like the random control within a
single episode ($\epsilon\approx 1$ over 2--3 acting rounds), so the value
of its Q-table is measured across episodes instead (below).

\subsection{Learning across episodes}

\todofig{Fig.~3: per-episode best over 20 episodes, shared Q-table/memory
vs.\ random control and static-TPE, two panels. Data:
\texttt{figs/episodes\_summary.csv} (FIGURES.md Fig.~3).}

With the Q-table and memory shared across 20 consecutive episodes,
\texttt{adapt-rl} improves as experience accumulates: on synthetic\_mlp the
mean per-episode best rises from 0.892 over the first five episodes to
0.949 over the last five (peak 0.975 at episode 5), while the unguided
control stays at 0.867$\to$0.855 and static TPE is flat at $\approx$0.75.
On Branin (minimize) \texttt{adapt-rl} improves from 0.734 to 0.524 (0.469
already at episode 5) while the unguided control drifts from 0.768 to
0.819; static TPE stays at $\approx$0.90. The Q-table stays small
(3 states on average), so learning remains tractable.

\subsection{Case study: detecting and repairing truncation}

\todofig{Fig.~4 (full width): (a) best-so-far with round actions;
(b) learning-rate interval per round (log scale) with GT 0.45 and sampled
points; (c) reward per round. Data: \texttt{figs/case\_study.json}
(FIGURES.md Fig.~4).}

Fig.~4 traces \texttt{adapt-rule-only} (seed 0) on synthetic\_mlp. After
round 0 (40 trials), cumulative importance marks learning\_rate as
dominant (0.876) and 75\% of the top-10\% trials sit at the \emph{upper}
boundary of the initial range $[10^{-3}, 10^{-1}]$:
\textsc{ShiftSearchRegion} moves the interval to
$[9.1{\times}10^{-3}, 9.1{\times}10^{-1}]$, preserving its (log-)width.
The warm-started TPE phase samples the newly opened region immediately:
the best jumps from 0.743 to 1.000, recovering learning rate 0.420
(ground truth 0.45), and a second \textsc{Shift} recenters the interval on
the good region. Round 2 spends the remaining budget and closes the
episode without a further action.

\subsection{Cost and reproducibility}

On rf\_digits, where each trial actually trains models, the
observe--decide--execute layer (diagnostics, rule engine, Q-updates,
memory, agent execution) costs 1.4\% of wall-clock time on average (1.7\%
worst case over the 60 adaptive runs); the remainder is model training.
On analytic objectives trials cost microseconds, so that layer trivially
dominates. The MIT-licensed toolkit ships 126 unit tests covering every
module; with fixed seeds, every table above regenerates from one command.

\section{Conclusion}
\label{sec:conclusion}

\toolname{} closes the loop from CAVE-style trajectory diagnosis to
schema-validated setup revision: guided adaptation substantially improves
results on misspecified setups at 1.4\% overhead, and the controller
improves across episodes. Future work: LLM backends in the loop and
embedding-based memory retrieval.

\bibliographystyle{unsrt}
\bibliography{refs}

\end{document}
